\documentclass[10pt]{article}

\usepackage[utf8]{inputenc}

\usepackage[T1]{fontenc}

\usepackage{amsmath}

\usepackage{amsfonts}

\usepackage{amssymb}

\usepackage[version=4]{mhchem}

\usepackage{stmaryrd}



\begin{document}

ное операторное обыкновенное диффференциальное уравнение на бесконечном промежутке



\[

\frac{d x}{d \lambda}=x(I-A x)(=(I-x A) x), x(0)=x_{0},

\]



где $A, x_{0}, I \in[X]$ ( $A, x_{0}$ - заданные, $I$ - тождественный оператор), $x(\lambda)$ - абстрактная функция со значениями из $[X](X-$ пространство типа $B)$. $К$ решению уравнения (2) приводится прямым методом вариации параметра задача построения обратного оператора $A^{-1}$ для обратимого оператора $A$, решение линейного операторного уравнения вида $I-A x=0$ и др. задачи. Если спектр оператора $A x_{0}\left(x_{0} A\right)$ расположен в правой полуплоскости, то задача (2) имеет единственное решение $x(\lambda), x(0)=x_{0}, \forall \lambda \in[0, \infty)$.



Аналитич. применение к задаче Коши (2) методов численного интегрирования обыкновенных дифференциальных уравнений приводит к целому классу прямых методов вариации параметра решения операторного уравнения (2), а следовательно, и ретения тех задач, к-рые сводятся к задаче (2). Напр., метод Эйлера с неравномерным шагом $h_{k}$ приводит к следующему методу построения оператора $A^{-1}$ :



\[

\begin{gathered}

x_{k+1}=x_{k}+h_{k} x_{k}\left(I-A x_{k}\right)\left(+h_{k}\left(I-x_{k} A\right) x_{k}\right), \\

k=0,1,2 \ldots

\end{gathered}

\]



Шаги $h_{k}$ в формуле (3) выбирают различными способами. Когда спектр оператора $I-A x_{0}\left(I-x_{0} A\right)$ расположен на действительной оси в интервале $\left[-\rho_{0}, 0\right]$, $0<\rho_{0}<\infty$, весьма әффективным является следующий способ выбора шагов $h_{k}$ :



\[

\begin{gathered}

h_{k}=1 /\left(1+\rho_{k}\right), \rho_{k+1}=\rho_{k} L_{k}, 4 L_{k}=\rho_{k} /\left(1+\rho_{k}\right), \\

k=0,1,2, \ldots, h_{k+1}=4 h_{k} /\left(1+h_{k}\right)^{2} .

\end{gathered}

\]



При этом быстрота сходимости метода (3) имеет порядок выше, чем $2^{k}$, норма вевязки убывает по закону $\rho_{k+1}=\rho_{k}^{2} / 4\left(1+\rho_{k}\right)$. Случай интервала $\left[0, \bar{\rho}_{0}\right], 0<\bar{\rho}_{0}<1$, выгодно сводится к рассмотренному одним шагом метода (3) при



\[

h=1 /\left(1-\bar{\rho}_{0}\right), \rho_{0}=\bar{\rho}_{0}^{2} / 4\left(1-\bar{\rho}_{0}\right) .

\]



Относительно сходимости метода (3), (4) для конкретных операторных уравнений можно высказать на основании общих фактов те или иные результаты в зависимости от того, какос иространство берется за основу.



К решению задачи (2) сводятся также задача построения проекторов вида



\[

P\left(A x_{0}\right)\left(P\left(x_{0} A\right)\right), P(K)=\left.e^{-K \lambda}\right|_{\lambda=\infty}, K \in[X],

\]



[случай, когда спектр огератора $I-A x_{0}\left(I-x_{0} A\right)$ pacположен в интервале $\left.\left[-\rho_{0}, 1\right]\right]$ и задача построения псевдообратных операторов $x^{+}(+x)$ для оператора $A$ таких, что



\[

I-A x^{+}=P\left(A x_{0}\right), I-{ }^{+} x A=P\left(x_{0} A\right), x^{+}={ }^{+} x .

\]



Для непосредственного построения проекторов $P\left(A x_{0}\right)$ $\left(P\left(x_{0} A\right)\right)$ метод (3), (4) можно перенисать в виде



\[

\begin{gathered}

P_{k+1}=P_{k}+h_{k} P_{k}\left(P_{k}-I\right), \quad k=0,1,2, \ldots, \\

P_{0}=I-A x_{0}\left(P_{0}=I-x_{0} A\right) .

\end{gathered}

\]



В зависимости от расположения спектра и свойств оператора $A$ в качестве $x_{0}$ выбирают, вапр., операторы $I, A, A^{*}$ или др.



К абстрактному линейному функциональному обыквовенному дифференциальному уравнению $(0 \leqslant \lambda \leqslant \infty)$



\[

\frac{d y}{d \lambda}=x_{0}(b-A y), y(0)=y_{0},

\]



где $b, y_{0} \in X, x_{0} . A \in[X], y(\lambda)-$ абстрактная функция со значевиями в $X$, сводится прямым методом вариации параметра задача непосредственного построения псевдорешений $y^{+}$линейных Ф. у. вида



\[

b-A y^{+}=P\left(A x_{0}\right)\left(b-A y_{0}\right)

\]



или Ф. у. $b-A y^{+}=0$, если $P\left(A x_{0}\right)\left(b-A y_{0}\right)=0$. К задаче (5) сводятся также и др. задачи, в том числе задачи для дифференциальных уравнений обыкновенных и с частными производными, интегральные уравнения и др. Формула



\[

y(\lambda)=e^{-x_{0} A \lambda} y_{0}+r(\lambda) x_{0} b, r(\lambda)=\int_{0}^{\lambda} e^{-x_{0} A(\lambda-s)} d s,

\]



дает единственное репеение уравнения (5), удовлетворяющее условию $y(0)=y_{0}$.



Применение к задаче (5), напр., метода Эйлера с шагом $h_{k+1}$ вперед приводит к следующему методу построения псевдорешений $y^{+} \Phi . \mathbf{y} .(6)$ :



\[

y_{k+1}=y_{k}+h_{k+1} x_{0}\left(b-A y_{k}\right), \quad k=0,1,2, \ldots

\]



Шаги $h_{k+1}$ выбирают, вапр., через корни многочленов Чебышева $T_{N}(t)$ :



\[

h_{j}=2\left[M+m-(M-m) t_{j}\right]^{-1}

\][



t\_\{j\}=\textbackslash cos \textbackslash left[\textbackslash left(2 x\_\{j\}-1\textbackslash right) \textbackslash pi / 2 N\textbackslash right], \textbackslash quad j=1,2, ···, N,

]



когда $x_{0} A$-самосопряженный оператор с ненулевой частью спектра на отрезке $[m, M], 0<m \leqslant M, P_{N}=$ $=\left(x_{1}, x_{2}, \ldots, x_{N}\right)-$ нек-рое упорядочение чисел $1,2, \ldots, N$ для устойчивости счета, $N$-порядок требуемого многотлена. Метод (7), (8) дает оптимальную оценку сходимости только ва $N$-м шаге. При $N=2^{i}$ (или $\approx 2 i$ ) весьма эффективным является выбор шагов в (7) последовательно через корни многочленов Чебышева



\[

\begin{gathered}

T_{2}(t)-T_{0}(t), T_{1}(t), T_{1}(t), T_{2}(t), T_{2}(t), \ldots \\

\ldots, T_{2^{i-2}}(t), T_{2^{i-2}}(t)

\end{gathered}

\]



вместо многочлена $T_{N}(t)$, что существенно упрощает задату упорядочения шаагов $h_{j}$ и повышает әффективность счета, особенно при больших $N$. В этом случае погрешность оптимально уменьшается после употребления всех упорядоченных корней каждого многочлена, что важно для упрощения контроля счета. В случае, если $P\left(A x_{0}\right)\left(b-A y_{0}\right)=0$, то $\quad P\left(x_{0} A\right) \times$ $\times\left(y^{+}-y_{0}\right)=0$; при этом, если $P\left(x_{0} A\right)-$ ортопроектор, то $y^{+}-y_{0}$ является нормальным решением Ф. у. $A\left(y^{+}-y_{0}\right)=b-A y_{0}$. С другой стороны, из (7) вытекает



$b-A y_{k+i}=U_{k+i}\left(b-A y_{0}\right), U_{k+1}=\prod_{i=1}^{k+1}\left(I-h_{i} A x_{0}\right)$,



\[

k=0,1,2, \ldots

\]



При әтом, если $P\left(A x_{0}\right)-$ проектор, то



\[

\left\|U_{k}-P\left(A x_{0}\right)\right\| \leqslant 2 q^{k} /\left(1+q^{2 k}\right) \rightarrow 0

\]



при $k \longrightarrow \infty(q=(\sqrt{M}-\sqrt{m}) /(\sqrt{M}+\sqrt{m}))$.



Существуют также прямые методы вариации параметра типа метода Эйлера, использующие рекуррентные соотнотения для многочленов Чебышева и близких к вим (без явного использования корней этих многочленов), с убыванием погрешности на каждом шаге. Эти методы являются многошаговыми и обладают повытевной быстротой сходимости. Применение к задаче (5) метода Эйлера с шагом $h_{k+1}$ назад дает следующий әффективный класс методов



\[

\begin{gathered}

y_{k+1}=y_{k}+h_{k+1}\left(I+h_{k+1} x_{0} A\right)^{-1} x_{0}\left(b-A y_{0}\right), \\

k=0,1,2, \ldots,

\end{gathered}

\]



откуда



\[

b-A y_{k+1}=W_{k+1}\left(b-A y_{0}\right),

\]



\[

W_{k+1}=\prod_{i=1}^{k+1}\left(I+h_{i} A x_{0}\right)^{-1}, \quad k=0,1,2, \ldots

\]





\end{document}